\section{Completion as Holonomy Elimination}
\label{sec:completion-holonomy}

Knuth--Bendix completion repairs a nonconfluent system by orienting nonjoinable critical pairs into new rules, and, in its general form, also simplifies or removes existing rules \cite{KnuthBendix1970,DershowitzJouannaud1990,BaaderNipkow1998}. We consider only the simplest kind of step, the addition of one rule. In the two rule-addition steps studied here, the overlap that produced the branching is still present after the new rule is added, so the branching remains; what changes is that it acquires a filler. We illustrate this with one string example and one term example, in each case adding a single rule by hand.

\begin{figure}[t]
\centering
\begin{tikzpicture}
  % string before
  \begin{scope}
    \node[st] (s) at (0,0) {$ab$};
    \node[nf] (x) at (-0.8,-1.1) {$x$};
    \node[nf] (y) at (0.8,-1.1) {$y$};
    \draw[rw] (s)--(x); \draw[rw] (s)--(y);
    \draw[defect] (x) .. controls (-0.3,-1.7) and (0.3,-1.7) .. (y);
    \node[font=\small] at (0,-2.3) {$E007$: before};
  \end{scope}
  \begin{scope}[xshift=2.9cm]
    \node[st] (s) at (0,0) {$ab$};
    \node[st] (x) at (-0.8,-1.1) {$x$};
    \node[nf] (y) at (0.8,-1.1) {$y$};
    \begin{scope}[on background layer]\fill[cellfill] (s.center)--(x.center)--(y.center)--cycle;\end{scope}
    \draw[rw] (s)--(x); \draw[rw] (s)--(y);
    \draw[rw, blue!60!black] (x)--node[below, font=\scriptsize]{new} (y);
    \node[font=\small] at (0,-2.3) {$E008$: after};
  \end{scope}
  % term before
  \begin{scope}[xshift=6.6cm]
    \node[st] (s) at (0,0) {$f(g(a))$};
    \node[nf] (l) at (-0.9,-1.1) {$f(b)$};
    \node[nf] (r) at (0.9,-1.1) {$p(a)$};
    \draw[rw] (s)--(l); \draw[rw] (s)--(r);
    \draw[defect] (l) .. controls (-0.3,-1.7) and (0.3,-1.7) .. (r);
    \node[font=\small] at (0,-2.3) {$TRS003$: before};
  \end{scope}
  \begin{scope}[xshift=9.9cm]
    \node[st] (s) at (0,0) {$f(g(a))$};
    \node[st] (l) at (-0.9,-1.1) {$f(b)$};
    \node[nf] (r) at (0.9,-1.1) {$p(a)$};
    \begin{scope}[on background layer]\fill[cellfill] (s.center)--(l.center)--(r.center)--cycle;\end{scope}
    \draw[rw] (s)--(l); \draw[rw] (s)--(r);
    \draw[rw, blue!60!black] (l)--node[below, font=\scriptsize]{new} (r);
    \node[font=\small] at (0,-2.3) {$TRS004$: after};
  \end{scope}
\end{tikzpicture}
\caption{Completion as holonomy elimination. Left: the string rules $ab \to x$ and $ab \to y$ produce an open branching with distinct normal forms $x$ and $y$; adding $x \to y$ fills it with join witness $y$. Right: the term rules $f(g(x)) \to p(x)$ and $g(a) \to b$ produce an open branching at $f(g(a))$; adding $f(b) \to p(a)$ fills it with join witness $p(a)$. In both cases the source and the two branches are unchanged.}
\label{fig:string-completion-before-after}
\label{fig:trs-completion-before-after}
\end{figure}

\subsection{The two examples}
\label{subsec:string-completion-holonomy}
\label{subsec:trs-completion-holonomy}

The string system $E007$ has the rules $ab \to x$ and $ab \to y$, whose left-hand sides coincide. From the start word $ab$ there is one branching, $x \leftarrow ab \to y$, and two normal forms, so the branching has nontrivial holonomy and the system is not confluent from $ab$. The system $E008$ adds the rule $x \to y$. The branching is still there, with the same source and branches, but now $x \to y$ provides a filler with join witness $y$. Its holonomy is trivial, and $y$ is the unique normal form.

The term system $TRS003$ has the rules $f(g(x)) \to p(x)$ and $g(a) \to b$, which overlap below the root. From the start term $f(g(a))$ there is one branching, $p(a) \leftarrow f(g(a)) \to f(b)$, whose branches are distinct normal forms. The system $TRS004$ adds the rule $f(b) \to p(a)$, which provides a filler with join witness $p(a)$.

\begin{table}[t]
    \caption{Before and after the added rule. The number of branchings does not change; the number of nonjoinable ones drops to zero, and confluence from the start term is restored.}
    \label{tab:holonomy-elimination-deltas}
    \centering
    \smallskip
    \small
\setlength{\tabcolsep}{5pt}
\begin{tabular}{@{}lllcccccc@{}}
\toprule
 & & & \multicolumn{2}{c}{Branchings} & \multicolumn{2}{c}{Nonjoinable} & \multicolumn{2}{c}{Confluent?} \\
\cmidrule(lr){4-5}\cmidrule(lr){6-7}\cmidrule(lr){8-9}
Domain & Before $\to$ after & Added rule & before & after & before & after & before & after \\
\midrule
string & $E007 \to E008$ & $x \to y$ & 1 & 1 & 1 & 0 & no & yes \\
term & $TRS003 \to TRS004$ & $f(b) \to p(a)$ & 1 & 1 & 1 & 0 & no & yes \\
\bottomrule
\end{tabular}

\end{table}

\begin{observation}[Completion as holonomy elimination]
\label{prop:completion-as-holonomy-elimination}
In the string pair $E007$--$E008$ and the term pair $TRS003$--$TRS004$, explored completely from their start terms, the added rule
\begin{enumerate}
    \item leaves the reachable critical-pair instance in place, with the same source and branches;
    \item changes its elementary holonomy from nontrivial to trivial;
    \item reduces the number of nonjoinable reachable instances from one to zero; and
    \item makes the system confluent from its start term.
\end{enumerate}
\end{observation}

In both pairs the ``after'' system consists of exactly the ``before'' rules plus the one added rule, with the same start term, so that the comparison is between a system and a one-rule extension of it. For the string pair this is checked automatically as part of the completion computation; for the term pair it can be read off the two rule sets.

\subsection{What is eliminated}
\label{subsec:what-is-eliminated}

The examples separate two quantities that are easily conflated: the number of generating branchings and the number of nontrivial ones. Completion reduces the second while leaving the first unchanged. In the reduction 2-complex, the effect of the new rule is to add a $1$-cell along which a $2$-cell can now be attached; the branching that was open becomes the boundary of a filled cell.

\begin{remark}[Scope]
\label{rem:completion-boundary}
The two examples are hand-built single steps, chosen so that one added rule suffices. They illustrate a reading of completion; they are not a geometric formulation of the completion procedure. In general, an added rule can create new critical pairs, completion may fail or fail to terminate, and orienting an equation requires a reduction ordering \cite{KnuthBendix1970,BaaderNipkow1998}. None of this is addressed here.
\end{remark}
